\documentclass[12pt,oneside]{book}

% ==================== METADATOS ====================
\newcommand{\universidad}{Universidad de Buenos Aires}
\newcommand{\facultad}{Facultad de Derecho}
\newcommand{\materia}{Concursos y Quiebras}
\newcommand{\titulo}{El concurso preventivo de las personas humanas no empresarias: ¿herramienta de protección frente al sobreendeudamiento o mecanismo susceptible de abuso?}
\newcommand{\autor}{Isaac Edgar Camacho Ocampo}
\newcommand{\anio}{2026}
\newcommand{\catedra}{Cátedra: Mag. Karina Bresca}
% TODO: contenido incompleto o ambiguo en las notas originales (no figura la carrera; el apunte deja sin completar el campo "Integrantes")

% ==================== PAQUETES ====================
\usepackage[utf8]{inputenc}
\usepackage[T1]{fontenc}
% La instalación del entorno LaTeX disponible aquí no incluye el soporte de babel en español,
% por lo que se usa inglés como idioma de compilación para asegurar que el índice se genere.
\usepackage[english]{babel}
\usepackage[top=2.5cm,bottom=2.5cm,left=3cm,right=2.5cm,headheight=15pt]{geometry}
\usepackage{amsmath}
\usepackage{amssymb}
\usepackage{mathtools}
\usepackage{graphicx}
\usepackage{float}
\usepackage{caption}
\usepackage{subcaption}
\usepackage{booktabs}
\usepackage{array}
\usepackage{longtable}
\usepackage{tabularx}
\usepackage{enumitem}
\usepackage{xcolor}
\usepackage[most]{tcolorbox}
\usepackage{fancyhdr}
\usepackage{ragged2e}
\usepackage[
    colorlinks=true,
    linkcolor=blue!50!black,
    citecolor=green!40!black,
    urlcolor=blue!60!black,
    pdftitle={\titulo},
    pdfauthor={\autor},
    pdfsubject={Apuntes universitarios},
    bookmarks=true
]{hyperref}

% ==================== CONFIGURACIÓN ====================
\graphicspath{{./img/}{./graficos/}{./figuras/}}
\setcounter{tocdepth}{2}
\setlength{\parindent}{0pt}
\setlength{\parskip}{0.5em}
\setlist[itemize]{topsep=0.3em,itemsep=0.2em}
\setlist[enumerate]{topsep=0.3em,itemsep=0.2em}
\clubpenalty=10000
\widowpenalty=10000

% ==================== COMANDOS ====================
\newcommand{\abs}[1]{\left|#1\right|}
\newcommand{\norm}[1]{\left\lVert#1\right\rVert}
\newcommand{\vect}[1]{\mathbf{#1}}
\newcommand{\deriv}[2]{\frac{d#1}{d#2}}
\newcommand{\pderiv}[2]{\frac{\partial#1}{\partial#2}}
\newcommand{\LCQ}{LCQ}

% ==================== ENTORNOS / CAJAS ====================
\newtcolorbox{definition}[1]{enhanced,breakable,title={Definición: #1},fonttitle=\bfseries,colback=blue!3,colframe=blue!50!black,boxrule=0.8pt,arc=2mm}
\newtcolorbox{important}{enhanced,breakable,title={Importante},fonttitle=\bfseries,colback=yellow!8,colframe=orange!70!black,boxrule=0.8pt,arc=2mm}
\newtcolorbox{example}[1]{enhanced,breakable,title={Ejemplo: #1},fonttitle=\bfseries,colback=green!4,colframe=green!40!black,boxrule=0.8pt,arc=2mm}
\newtcolorbox{formula}[1]{enhanced,breakable,title={Fórmula / Regla: #1},fonttitle=\bfseries,colback=purple!4,colframe=purple!50!black,boxrule=0.8pt,arc=2mm}
\newtcolorbox{exercise}[1]{enhanced,breakable,title={Ejercicio: #1},fonttitle=\bfseries,colback=gray!5,colframe=gray!60!black,boxrule=0.8pt,arc=2mm}
\newtcolorbox{note}{enhanced,breakable,title={Nota},fonttitle=\bfseries,colback=white,colframe=gray!50,boxrule=0.6pt,arc=2mm}

% Columnas del cuadro Cornell
\newcolumntype{Q}{>{\raggedright\arraybackslash\bfseries}p{0.26\textwidth}}
\newcolumntype{R}{>{\raggedright\arraybackslash}p{0.66\textwidth}}

% ==================== CONFIGURACIÓN FINAL ====================
\pagestyle{fancy}
\fancyhf{}
\fancyhead[L]{\nouppercase{\leftmark}}
\fancyhead[R]{\materia}
\fancyfoot[C]{\thepage}
\renewcommand{\headrulewidth}{0.4pt}
\renewcommand{\footrulewidth}{0pt}
\fancypagestyle{plain}{
    \fancyhf{}
    \fancyfoot[C]{\thepage}
    \renewcommand{\headrulewidth}{0pt}
}

\begin{document}

% ==================== PORTADA ====================
\begin{titlepage}
\thispagestyle{empty}
\begin{center}
\vspace*{1cm}
{\large \textsc{\universidad}}\\[0.2cm]
{\large \textsc{\facultad}}\\[3cm]
{\Huge \textbf{\materia}}\\[1cm]
{\Large \titulo}\\[0.8cm]
\textit{Comprender es el primer paso para convertir el estudio en conocimiento.}\\[1.2cm]
\rule{0.70\textwidth}{0.5pt}\\[0.5cm]
{\large Apuntes de estudio}\\[0.5cm]
\rule{0.70\textwidth}{0.5pt}\\[0.8cm]
{\normalsize \catedra}
\vfill
{\large \autor}\\[0.3cm]
{\large \anio}
\end{center}
\vfill
\begin{flushleft}
\textit{Este es un modesto aporte a los estudiantes de ``Concursos y Quiebras'' no pretende sustituir el material oficial de la materia.}
\end{flushleft}
\end{titlepage}

% ==================== ÍNDICE ====================
\frontmatter
\tableofcontents
\mainmatter

% ==================== CONTENIDO ====================

% -------------------- A --------------------
\chapter[Introducción y planteo del problema]{Introducción y planteo del problema}

Este material aborda una problemática de creciente relevancia para el derecho concursal argentino: la aplicación del régimen de concurso preventivo previsto por la Ley 24.522 a las personas humanas que no desarrollan una actividad empresarial o comercial organizada y que atraviesan una situación de sobreendeudamiento.

El problema surge de una particularidad del sistema argentino. La Ley de Concursos y Quiebras (LCQ) admite que una persona humana solicite su propio concurso preventivo, aunque no sea comerciante, empresario ni titular de una organización productiva. Sin embargo, esa amplitud en cuanto a los sujetos que pueden acceder al régimen no estuvo acompañada por la creación de un procedimiento específicamente diseñado para atender la insolvencia de los consumidores, trabajadores, jubilados o profesionales que carecen de una estructura empresarial.

En la actualidad, una persona puede contraer deudas mediante tarjetas de crédito, préstamos personales, financiación de bienes de consumo o plataformas digitales. En muchos casos, el endeudamiento se origina en la necesidad de afrontar gastos ordinarios, sostener el hogar o cubrir una contingencia imprevista. La pérdida del empleo, una enfermedad, una disminución de los ingresos o el aumento del costo de vida pueden transformar obligaciones inicialmente asumibles en una carga económica imposible de afrontar.

El sobreendeudamiento de una persona humana no empresaria, por lo tanto, no responde necesariamente a una conducta especulativa o irresponsable. Puede ser consecuencia de circunstancias ajenas a su voluntad, de una planificación económica deficiente, de una acumulación progresiva de créditos o de una combinación de esos factores. También puede presentarse, desde luego, en situaciones en las que el deudor haya actuado de mala fe o pretenda utilizar el proceso concursal para perjudicar a sus acreedores.

Esta diversidad de situaciones plantea una dificultad jurídica: \textbf{¿cómo debe responder el derecho concursal frente a la insolvencia de una persona humana que no desarrolla actividad empresarial, sin convertir el concurso preventivo en un beneficio indiscriminado para cualquier deudor ni dejar sin protección a quien realmente necesita una solución?}

La Ley 24.522 establece un procedimiento colectivo cuya finalidad es ordenar las relaciones entre el deudor y sus acreedores, evitar la dispersión de las ejecuciones individuales y procurar una solución preventiva que permita superar la cesación de pagos. No obstante, buena parte de sus exigencias y etapas responde a una lógica empresarial: documentación contable, intervención sindical, publicación de edictos, verificación de créditos, categorización de acreedores y negociación de un acuerdo preventivo.

Si bien el régimen contempla una modalidad simplificada para los pequeños concursos, regulada en los artículos 288 y 289, esa simplificación no configura un procedimiento autónomo para el consumidor sobreendeudado. En consecuencia, una persona que percibe un salario y tiene deudas derivadas de préstamos personales puede quedar sometida, en términos generales, a un proceso estructurado sobre parámetros semejantes a los aplicables a una empresa.

\begin{important}
\textbf{Pregunta de investigación.} ¿El régimen de concurso preventivo previsto por la Ley 24.522 constituye una herramienta adecuada para proteger a las personas humanas no empresarias afectadas por el sobreendeudamiento, sin desatender los derechos de los acreedores ni facilitar la utilización abusiva del proceso concursal?
\end{important}

La relevancia de esta pregunta es jurídica y social. Por un lado, el deudor de buena fe necesita una vía efectiva para ordenar sus obligaciones, evitar que la multiplicidad de ejecuciones individuales agrave su situación y recuperar, en la medida de lo posible, su capacidad económica. Por otro lado, los acreedores tienen derecho a que sus créditos sean reconocidos y satisfechos conforme a las reglas legales, sin quedar expuestos a maniobras dilatorias, ocultamientos patrimoniales o propuestas de acuerdo abusivas.

\begin{important}
\textbf{Hipótesis central.} El régimen argentino permite el acceso de las personas humanas al concurso preventivo, pero no ofrece una respuesta suficientemente adaptada a la insolvencia del consumidor no empresario. Las exigencias formales, los costos del procedimiento y la ausencia de un mecanismo específico de segunda oportunidad pueden dificultar el acceso de los deudores de buena fe. Al mismo tiempo, la flexibilidad de las propuestas de acuerdo y la falta de un procedimiento diferenciado exigen controles judiciales adecuados para evitar el abuso.
\end{important}

La posición sostenida no consiste en afirmar que toda persona sobreendeudada deba recibir automáticamente una liberación de sus deudas, ni que el concurso preventivo deba ser restringido a quienes desarrollan una empresa. Se considera que debe existir una solución equilibrada: \textbf{protección efectiva para el deudor de buena fe y tutela suficiente para los acreedores frente a conductas abusivas}.

Para desarrollar esta hipótesis se analizan las disposiciones de la Ley 24.522, las posiciones doctrinarias vinculadas con la insolvencia de las personas humanas, los precedentes judiciales relacionados con el acceso al concurso y el control de los acuerdos preventivos, y algunas soluciones del derecho comparado.

\section*{Cuadro Cornell}
\addcontentsline{toc}{section}{Cuadro Cornell}
\begin{longtable}{QR}
\toprule
Preguntas / palabras clave & \textbf{Notas / respuestas} \\
\midrule
\endhead
¿Qué problema aborda el trabajo? &
Cómo debe responder el derecho concursal a la insolvencia de una persona humana que no desarrolla actividad empresarial, sin convertir el concurso preventivo en un beneficio indiscriminado ni dejar sin protección a quien necesita una solución. \\[0.8em]
¿Qué particularidad tiene el sistema argentino? &
La LCQ admite que una persona humana solicite su propio concurso preventivo aunque no sea comerciante ni empresario, pero no creó un procedimiento específicamente diseñado para la insolvencia de consumidores, trabajadores, jubilados o profesionales sin estructura empresarial. \\[0.8em]
¿Por qué puede no ser adecuado el procedimiento para el consumidor? &
Porque buena parte de las exigencias y etapas de la LCQ responde a una lógica empresarial (documentación contable, intervención sindical, edictos, verificación de créditos, categorización de acreedores y negociación del acuerdo). Los pequeños concursos (arts.\ 288 y 289) simplifican, pero no configuran un procedimiento autónomo para el consumidor sobreendeudado. \\[0.8em]
¿Cuál es la hipótesis central? &
El régimen permite el acceso de las personas humanas, pero no ofrece una respuesta suficientemente adaptada a la insolvencia del consumidor no empresario; las exigencias formales, los costos y la ausencia de segunda oportunidad dificultan el acceso del deudor de buena fe, y la flexibilidad de las propuestas exige controles judiciales para evitar el abuso. \\[0.8em]
¿Cuál es la posición sostenida? &
Una solución equilibrada: protección efectiva para el deudor de buena fe y tutela suficiente para los acreedores frente a conductas abusivas. \\
\midrule
\multicolumn{2}{p{\dimexpr0.92\textwidth+2\tabcolsep\relax}}{\textbf{Resumen:} El trabajo examina si el concurso preventivo de la Ley 24.522 protege adecuadamente a las personas humanas no empresarias sobreendeudadas sin desatender a los acreedores ni facilitar el abuso. Su hipótesis es que el acceso existe pero la respuesta no está adaptada, y su posición es una solución equilibrada.} \\
\bottomrule
\end{longtable}

% -------------------- B --------------------
\chapter[Análisis normativo]{Análisis normativo}

\section{El presupuesto subjetivo: la persona humana como sujeto concursable}

El artículo 2 de la Ley 24.522 establece quiénes pueden ser declarados en concurso. La norma adopta un criterio amplio e incluye a las personas humanas ---en la terminología original de la ley, personas de existencia visible---, sin exigir que desarrollen una actividad comercial o empresarial.

Esta disposición representa un cambio importante respecto de los sistemas que vinculaban el acceso a los procedimientos concursales con la condición de comerciante. Bajo el régimen vigente, la posibilidad de solicitar el concurso preventivo no depende de que el deudor tenga una empresa, lleve una contabilidad comercial o participe de una actividad económica organizada.

Por lo tanto, pueden solicitar su concurso preventivo, si reúnen los presupuestos legales, un trabajador en relación de dependencia, un jubilado, un profesional independiente o un consumidor que haya contraído obligaciones que no puede afrontar regularmente.

El artículo 5 de la LCQ reconoce la legitimación para solicitar el concurso preventivo a los sujetos comprendidos en el artículo 2. A su vez, el artículo 10 establece que el concurso preventivo puede ser solicitado mientras la quiebra no haya sido declarada.

De esta manera, el ordenamiento reconoce que la insolvencia no es un fenómeno exclusivo de las empresas. También puede afectar a las personas humanas que no desarrollan actividad productiva organizada.

\begin{important}
\textbf{Que una persona sea sujeto concursable no significa que pueda acceder al concurso por el solo hecho de encontrarse endeudada.} Para ello debe reunir, además, el presupuesto objetivo exigido por la ley.
\end{important}

\section{El presupuesto objetivo: cesación de pagos y sobreendeudamiento}

El artículo 1 de la LCQ establece que el estado de cesación de pagos, cualquiera sea su causa y la naturaleza de las obligaciones afectadas, constituye el presupuesto objetivo para la apertura de los concursos regulados por la ley, con las excepciones que ella contempla.

\begin{definition}{Cesación de pagos}
Estado patrimonial general y permanente de imposibilidad de cumplir regularmente las obligaciones exigibles mediante recursos ordinarios. No debe confundirse con un incumplimiento aislado, una dificultad momentánea de liquidez o el simple hecho de tener un nivel elevado de deuda.
\end{definition}

El artículo 78 dispone que el estado de cesación de pagos debe exteriorizarse mediante hechos reveladores. El artículo 79 enumera, entre otros, algunos de esos hechos:
\begin{itemize}
    \item reconocimiento judicial o extrajudicial del estado de cesación de pagos;
    \item mora en el cumplimiento de una obligación;
    \item ocultación o ausencia del deudor;
    \item clausura de establecimientos;
    \item venta a precio vil de bienes;
    \item revocación judicial de actos realizados en fraude de los acreedores;
    \item otros supuestos que permitan advertir la imposibilidad de cumplir regularmente.
\end{itemize}

La doctrina ha elaborado distintas teorías para explicar la cesación de pagos:
\begin{itemize}
    \item \textbf{Concepción restringida:} la vinculaba con el incumplimiento de una obligación.
    \item \textbf{Concepción intermedia:} exigía una pluralidad de incumplimientos.
    \item \textbf{Concepción amplia o moderna:} la identifica con un estado patrimonial de impotencia para afrontar regularmente las obligaciones exigibles, que puede manifestarse a través de distintos hechos reveladores.
\end{itemize}

Esta última concepción resulta especialmente importante para analizar la insolvencia de las personas humanas. Un consumidor puede continuar pagando algunas obligaciones mediante nuevos préstamos, refinanciaciones o utilización de tarjetas de crédito, aunque su situación patrimonial ya presente un desequilibrio grave. La apariencia de cumplimiento no necesariamente significa que tenga capacidad económica real para sostener sus compromisos.

No obstante, es necesario mantener una distinción conceptual: \textbf{el sobreendeudamiento no equivale automáticamente a la cesación de pagos}. Una persona puede tener muchas deudas y todavía contar con ingresos o bienes suficientes para cumplirlas regularmente. A la inversa, puede existir cesación de pagos aun cuando el monto total de las obligaciones no sea extraordinariamente elevado, si el deudor carece de medios regulares para afrontarlas.

El problema identificado es que el sistema concursal argentino exige, como regla general, que la insolvencia haya alcanzado el estado de cesación de pagos. No contempla un procedimiento preventivo específico para la persona humana que, sin encontrarse todavía en ese estado, presenta signos objetivos de sobreendeudamiento y necesita reorganizar sus obligaciones antes de que la situación se agrave.

Por ello, se considera que una eventual reforma debería estudiar la incorporación de mecanismos de prevención o alerta temprana, sin confundirlos con el presupuesto objetivo que actualmente exige el artículo 1 de la LCQ.

\section{Los requisitos de presentación: artículo 11 de la LCQ}

El artículo 11 de la Ley 24.522 establece los requisitos formales que debe cumplir el deudor que solicita la apertura de su concurso preventivo. Entre ellos se encuentran la explicación de las causas concretas de la situación patrimonial, la indicación de la época en que se produjo la cesación de pagos, la presentación de un estado detallado y valorado del activo y del pasivo, la documentación relativa a los acreedores y, cuando corresponda, la información contable y laboral exigida por la norma.

Estas exigencias cumplen una finalidad importante: permitir que el juez, el síndico y los acreedores conozcan la situación patrimonial del deudor y puedan controlar la información que se presenta.

Sin embargo, aparece una dificultad cuando el solicitante es una persona humana que no desarrolla una actividad empresarial. Un trabajador asalariado o un jubilado normalmente no lleva libros de comercio, no confecciona balances ni posee una contabilidad organizada semejante a la de una sociedad comercial.

El problema no consiste en que el deudor esté exento de informar su situación patrimonial. Esa información resulta indispensable. La dificultad radica en exigirle instrumentos contables que pueden no corresponder a su actividad o que resultan desproporcionados para la naturaleza de sus ingresos y obligaciones.

La propia LCQ contiene una respuesta parcial a esta dificultad en el régimen de pequeños concursos. El artículo 289 exime, en los casos comprendidos en ese régimen, de los requisitos previstos por el artículo 11, incisos 3 y 5. Sin embargo, esa simplificación no elimina todas las exigencias ni crea un procedimiento especial para la insolvencia del consumidor.

Por eso, se entiende que la información patrimonial debe continuar siendo un requisito esencial, pero debería adecuarse a la realidad del sujeto concursado. En el caso de una persona humana no empresaria, podrían resultar más apropiados los recibos de sueldo, las constancias de ingresos, los resúmenes bancarios, los estados de cuenta, los contratos de préstamo, los resúmenes de tarjetas y una declaración patrimonial suficientemente detallada.

La finalidad no debe ser reducir el control judicial, sino permitir que ese control se realice mediante documentación adecuada a la situación del deudor.

\section{Los pequeños concursos: artículos 288 y 289}

Los artículos 288 y 289 de la LCQ regulan los pequeños concursos y quiebras. Este régimen constituye la principal simplificación que la legislación vigente contempla para los procesos concursales de menor dimensión económica.

Conforme al artículo 288, se consideran pequeños concursos aquellos en los que se verifique indistintamente cualquiera de las siguientes circunstancias:
\begin{itemize}
    \item que el pasivo denunciado no alcance el equivalente a trescientos salarios mínimos, vitales y móviles;
    \item que el proceso no presente más de veinte acreedores quirografarios;
    \item que el deudor no posea más de veinte trabajadores en relación de dependencia.
\end{itemize}

La utilización de la expresión \textbf{``indistintamente''} significa que no es necesario que se reúnan simultáneamente los tres requisitos. La configuración de cualquiera de ellos permite encuadrar el proceso como pequeño concurso.

El artículo 289 establece determinadas simplificaciones:
\begin{itemize}
    \item eximición de los requisitos contables de los incisos 3 y 5 del artículo 11;
    \item ausencia de constitución de los comités de acreedores;
    \item inaplicabilidad del procedimiento de salvataje previsto en el artículo 48;
    \item control del cumplimiento del acuerdo a cargo del síndico.
\end{itemize}

La existencia de este régimen demuestra que el legislador reconoció la necesidad de simplificar determinados procedimientos concursales. Sin embargo, la solución es limitada:
\begin{enumerate}
    \item Los criterios del artículo 288 son cuantitativos. No distinguen necesariamente entre una pequeña empresa y una persona humana que se endeudó para afrontar gastos de subsistencia. Ambos sujetos pueden quedar comprendidos en el régimen, aunque sus realidades económicas y sus necesidades sean diferentes.
    \item Las simplificaciones del artículo 289 no eliminan las principales etapas del concurso preventivo. El deudor continúa sujeto a la intervención de la sindicatura, a la publicidad concursal, al período de verificación de créditos, a la negociación del acuerdo y a las mayorías legales.
\end{enumerate}

\begin{important}
\textbf{El pequeño concurso no equivale a un procedimiento especial de insolvencia del consumidor.} Es una modalidad simplificada del régimen general, pero no ofrece una respuesta integral a las particularidades de la persona humana no empresaria.
\end{important}

Este aspecto se relaciona con una observación formulada en el marco de la cursada: una persona que perdió su empleo y no puede afrontar sus deudas puede quedar sometida a un procedimiento prácticamente igual al de una organización empresarial, pese a que sus posibilidades de generar ingresos, su estructura patrimonial y la naturaleza de sus obligaciones son muy diferentes.

\section{Publicidad concursal y costos de acceso: artículos 27 y 30}

Los artículos 27 y 30 de la LCQ regulan aspectos relevantes de la publicidad del concurso preventivo y las consecuencias del incumplimiento de las cargas impuestas al concursado.

La publicación de edictos cumple una función esencial: dar publicidad a la apertura del proceso y permitir que los acreedores y terceros conozcan la existencia del concurso. No se trata de una formalidad carente de finalidad. La publicidad contribuye a la universalidad del proceso, a la transparencia y a la posibilidad de que todos los interesados ejerzan sus derechos.

Sin embargo, en el caso de una persona humana sobreendeudada, el costo de cumplir con las publicaciones puede convertirse en una dificultad importante. El deudor que solicita la apertura de su concurso precisamente porque no puede afrontar sus obligaciones puede carecer de recursos para solventar los gastos de publicidad.

La situación plantea una tensión entre dos intereses legítimos: por un lado, la necesidad de garantizar la publicidad del proceso; por otro, el derecho de acceso a la justicia y la posibilidad real de que el deudor pueda sostener el trámite.

No se considera correcto eliminar la publicidad concursal, porque ello podría perjudicar a los acreedores que no fueron denunciados o que necesitan tomar conocimiento del proceso. Pero sí se entiende que debe analizarse si los medios de publicidad previstos por la legislación resultan adecuados para todos los tipos de deudores.

Una alternativa sería ampliar los mecanismos digitales de publicidad oficial y permitir formas complementarias de notificación a los acreedores identificados, siempre que se preserve la posibilidad de conocimiento por parte de terceros y se respeten las garantías del proceso colectivo.

\section{El control del deudor y de la propuesta: artículos 15, 16, 17, 43 y 52}

La preocupación por el posible uso abusivo del concurso preventivo exige analizar las facultades de control que contempla la LCQ.

Desde la apertura del concurso, el deudor conserva la administración de su patrimonio, pero lo hace bajo vigilancia de la sindicatura, conforme al artículo 15. Los artículos 16 y 17 establecen límites a determinados actos de disposición y regulan las consecuencias de su incumplimiento. Estas disposiciones muestran que la apertura del concurso no significa que el deudor pueda disponer libremente de sus bienes ni actuar sin controles. La ley procura evitar que, durante el proceso, se altere injustificadamente la situación patrimonial en perjuicio de los acreedores.

Por otra parte, el artículo 43 regula el contenido de la propuesta de acuerdo preventivo. El deudor puede ofrecer distintas alternativas de reorganización, entre ellas quitas, esperas u otras modalidades admitidas por la legislación. La amplitud de las propuestas permite adaptar el acuerdo a la realidad económica de cada deudor. Sin embargo, también genera el riesgo de que se formulen propuestas que, aunque obtengan las mayorías legales, resulten abusivas o contrarias a la finalidad del proceso.

\begin{important}
\textbf{Artículo 52, inciso 4, de la LCQ:} el juez no homologará una propuesta abusiva o en fraude a la ley. La obtención de las mayorías legales no significa que cualquier propuesta deba ser homologada automáticamente: el juez conserva la obligación de controlar que el acuerdo no resulte abusivo ni fraudulento.
\end{important}

Esta norma constituye un límite al ejercicio de la autonomía del deudor y de los acreedores. Este control debe interpretarse en conjunto con los artículos 9 y 10 del Código Civil y Comercial de la Nación, que consagran los principios de buena fe y prohibición del abuso del derecho.

También debe considerarse el artículo 42 de la Constitución Nacional y la Ley 24.240 de Defensa del Consumidor, cuando la situación se origine en relaciones de consumo. En particular, el artículo 36 de la Ley 24.240 contiene previsiones específicas para las operaciones financieras y de crédito para el consumo.

De esta manera, el análisis del sobreendeudamiento no puede concentrarse exclusivamente en la conducta del deudor. También debe considerar la actuación de quienes conceden el crédito, especialmente cuando se trata de proveedores profesionales que cuentan con mayores herramientas para evaluar la solvencia y la capacidad de pago del consumidor.

Esto no significa que el acreedor sea responsable automáticamente por el incumplimiento. Significa que la evaluación de la buena fe y de las circunstancias del endeudamiento debe considerar el comportamiento de todos los sujetos involucrados.

\section{El derecho comparado: la exoneración del pasivo insatisfecho en España}

El derecho comparado ofrece herramientas útiles para analizar posibles alternativas al régimen argentino. En particular, resulta relevante el sistema español de segunda oportunidad y la denominada \textbf{Exoneración del Pasivo Insatisfecho (EPI)}.

La legislación española, reformada mediante la Ley 16/2022, de 5 de septiembre, en el marco de la adaptación a la Directiva (UE) 2019/1023, contempla mecanismos que permiten a determinadas personas físicas insolventes obtener la exoneración de deudas que no pueden ser satisfechas, siempre que cumplan los requisitos legales. En términos generales, el sistema contempla vías que pueden vincularse con la liquidación del patrimonio o con un plan de pagos, bajo determinadas condiciones y controles.

La finalidad es permitir que el deudor de buena fe pueda superar una situación de insolvencia sin permanecer indefinidamente condicionado por obligaciones que, en la práctica, no tiene posibilidades razonables de satisfacer.

\begin{note}
La EPI española no equivale al concurso preventivo argentino ni debe confundirse con la rehabilitación del fallido regulada por la LCQ. Se trata de institutos diferentes, pertenecientes a sistemas normativos distintos.
\end{note}

La comparación resulta útil porque demuestra que existen modelos que procuran equilibrar dos objetivos: permitir una segunda oportunidad al deudor honesto y establecer controles para excluir o limitar el beneficio cuando existen conductas fraudulentas o contrarias a la buena fe.

No se propone trasladar automáticamente el sistema español al derecho argentino. Se considera, sin embargo, que constituye una referencia valiosa para discutir la necesidad de incorporar un régimen específico de insolvencia de las personas humanas no empresarias.

\section*{Cuadro Cornell}
\addcontentsline{toc}{section}{Cuadro Cornell}
\begin{longtable}{QR}
\toprule
Preguntas / palabras clave & \textbf{Notas / respuestas} \\
\midrule
\endhead
¿Quiénes pueden ser declarados en concurso según el art.\ 2 de la LCQ? &
Las personas humanas, sin que se exija actividad comercial o empresarial. Pueden solicitarlo, si reúnen los presupuestos legales, un trabajador en relación de dependencia, un jubilado, un profesional independiente o un consumidor. El art.\ 5 reconoce la legitimación a los sujetos del art.\ 2 y el art.\ 10 permite solicitarlo mientras la quiebra no haya sido declarada. \\[0.8em]
¿Basta con ser sujeto concursable para acceder al concurso? &
No. Debe reunirse además el presupuesto objetivo exigido por la ley. \\[0.8em]
¿Qué es la cesación de pagos y dónde se establece como presupuesto objetivo? &
Es un estado patrimonial general y permanente de imposibilidad de cumplir regularmente las obligaciones exigibles mediante recursos ordinarios. El art.\ 1 la establece como presupuesto objetivo; no se confunde con un incumplimiento aislado, una dificultad momentánea de liquidez ni con un nivel elevado de deuda. \\[0.8em]
¿Cómo se exterioriza la cesación de pagos y qué teorías existen? &
Mediante hechos reveladores (arts.\ 78 y 79). La concepción restringida la vinculaba con el incumplimiento de una obligación; la intermedia, con una pluralidad de incumplimientos; la amplia o moderna, con un estado patrimonial de impotencia para afrontar regularmente las obligaciones exigibles. \\[0.8em]
¿Es lo mismo sobreendeudamiento que cesación de pagos? &
No. Puede haber muchas deudas con ingresos o bienes suficientes para cumplirlas, y puede haber cesación de pagos sin un monto total extraordinario si el deudor carece de medios regulares. El sistema exige la cesación y no contempla un procedimiento preventivo específico para el sobreendeudamiento previo a ese estado. \\[0.8em]
¿Qué dificultad plantea el art.\ 11 para la persona humana no empresaria? &
Exige requisitos formales (causas, época de la cesación, estado detallado del activo y pasivo, documentación sobre acreedores, información contable y laboral). El deudor no está exento de informar su situación patrimonial, pero se le pueden exigir instrumentos contables desproporcionados. Se proponen recibos de sueldo, constancias de ingresos, resúmenes bancarios y de tarjetas, contratos de préstamo y declaración patrimonial detallada. \\[0.8em]
¿Qué son los pequeños concursos y por qué no resuelven el problema? &
Regulados por los arts.\ 288 y 289; se configuran con cualquiera de estas circunstancias: pasivo menor a trescientos salarios mínimos, vitales y móviles; hasta veinte acreedores quirografarios; hasta veinte trabajadores. El art.\ 289 simplifica, pero no elimina las etapas principales ni crea un procedimiento especial para el consumidor; sus criterios son cuantitativos. \\[0.8em]
¿Qué tensión plantean los edictos? &
Entre la publicidad del proceso (universalidad, transparencia, derechos de acreedores y terceros) y el acceso a la justicia del deudor que puede carecer de recursos para costearlos. Se propone analizar mecanismos digitales de publicidad oficial y notificaciones complementarias a los acreedores identificados. \\[0.8em]
¿Qué controles existen frente al abuso? &
Administración del patrimonio bajo vigilancia de la sindicatura (art.\ 15), límites a actos de disposición (arts.\ 16 y 17), y el art.\ 52 inc.\ 4, que impide homologar propuestas abusivas o en fraude a la ley, en conjunto con los arts.\ 9 y 10 del CCCN. \\[0.8em]
¿Qué es la EPI española y qué aporta a la comparación? &
Exoneración del Pasivo Insatisfecho: mecanismo de segunda oportunidad (Ley 16/2022, adaptación de la Directiva (UE) 2019/1023) que permite a determinadas personas físicas insolventes obtener la exoneración de deudas insatisfechas bajo requisitos y controles. No equivale al concurso preventivo argentino ni a la rehabilitación del fallido; es una referencia para discutir un régimen específico, sin trasladarlo automáticamente. \\
\midrule
\multicolumn{2}{p{\dimexpr0.92\textwidth+2\tabcolsep\relax}}{\textbf{Resumen:} La LCQ admite el concurso de cualquier persona humana (art.\ 2), pero exige cesación de pagos (art.\ 1), distinta del mero sobreendeudamiento. Los requisitos del art.\ 11, los pequeños concursos (arts.\ 288 y 289) y la publicidad concursal (arts.\ 27 y 30) no están pensados para el consumidor, aunque existen controles contra el abuso (arts.\ 15, 16, 17, 43 y 52). La EPI española es una referencia comparada.} \\
\bottomrule
\end{longtable}

% -------------------- C --------------------
\chapter[Doctrina]{Doctrina}

\section{La posición protectoria: la necesidad de una respuesta para la persona humana sobreendeudada}

Una primera línea de análisis doctrinario sostiene que el derecho concursal debe dar una respuesta adecuada a la insolvencia de las personas humanas y que no puede limitarse a reproducir las soluciones concebidas para las empresas.

Desde esta perspectiva, el sobreendeudamiento del consumidor constituye un problema que excede el incumplimiento individual de una obligación. Puede comprometer la subsistencia del deudor y de su grupo familiar, afectar su capacidad de acceder a bienes y servicios esenciales y provocar su exclusión económica.

Esta postura destaca que el endeudamiento puede originarse en circunstancias muy diferentes. No es lo mismo una persona que toma créditos para sostener gastos básicos después de perder su empleo que otra que oculta deliberadamente su situación patrimonial o contrae obligaciones sin intención de cumplirlas. Por ello, la respuesta jurídica no debería ser idéntica en todos los casos.

La posición protectoria también pone el acento en la conducta de los proveedores de crédito. Las entidades financieras y las empresas que ofrecen préstamos al consumo tienen una posición técnica y económica diferente de la del consumidor. Cuentan con información, sistemas de evaluación y herramientas de análisis que pueden permitirles conocer la capacidad de pago del solicitante.

\begin{definition}{Crédito responsable}
Idea según la cual el otorgamiento de financiación debe considerar la capacidad real de repago y evitar prácticas que favorezcan la acumulación de obligaciones que el consumidor difícilmente podrá afrontar.
\end{definition}

Esta posición no implica sostener que toda deuda de consumo deba ser perdonada ni que el acreedor pierda su derecho a cobrar. Su fundamento es que la insolvencia de la persona humana debe ser abordada mediante mecanismos que permitan una solución razonable y eviten la exclusión económica permanente.

En esta línea de debate se ubican los trabajos doctrinarios sobre el sobreendeudamiento del consumidor y la necesidad de revisar el tratamiento concursal de la persona humana, entre ellos los desarrollos atribuidos en las versiones preliminares a Junyent Bas y Molina Sandoval, así como los estudios sobre procedimientos concursales y tutela de la persona humana.

\section{La posición de prevención del abuso: tutela del crédito y seguridad jurídica}

Una segunda perspectiva advierte sobre los riesgos de flexibilizar excesivamente el régimen concursal de las personas humanas.

Desde esta posición, el concurso preventivo no puede convertirse en un mecanismo que permita al deudor liberarse de sus obligaciones mediante propuestas irrazonables o utilizar el proceso únicamente para paralizar las acciones individuales de cobro.

El crédito constituye un elemento central de las relaciones económicas. Las personas y las empresas conceden financiación sobre la base de la expectativa legítima de que las obligaciones serán cumplidas. Si el sistema permite que cualquier deudor obtenga una reducción sustancial de sus obligaciones sin controles adecuados, puede debilitarse la confianza en las relaciones crediticias.

Esta postura destaca, además, la importancia de la igualdad entre los acreedores. El proceso concursal debe evitar que algunos acreedores obtengan ventajas indebidas en perjuicio de otros y debe impedir que el deudor altere fraudulentamente su patrimonio antes o durante el procedimiento.

El riesgo de abuso puede manifestarse de diferentes maneras:
\begin{itemize}
    \item ocultamiento de bienes;
    \item simulación de deudas;
    \item contratación de nuevos préstamos sin intención de pago;
    \item utilización del concurso con fines meramente dilatorios;
    \item presentación de propuestas que no guarden una relación razonable con las posibilidades económicas del deudor.
\end{itemize}

\begin{note}
No corresponde atribuir a toda la doctrina tradicional una oposición absoluta al concurso de las personas humanas. Autores como Rivera y Rouillon desarrollan el régimen general de la LCQ y sus mecanismos de protección del crédito, pero no corresponde afirmar, sin una cita específica, que sostengan una posición contraria a toda tutela concursal del consumidor.
\end{note}

La preocupación por el abuso no conduce necesariamente a negar el acceso al concurso. Conduce a exigir controles adecuados y a preservar los derechos de los acreedores.

\section{Comparación de las posiciones}

Las dos perspectivas parten de preocupaciones legítimas, pero colocan el acento en aspectos diferentes.

\begin{table}[H]
\centering
\renewcommand{\arraystretch}{1.3}
\begin{tabularx}{\textwidth}{>{\raggedright\arraybackslash\bfseries}p{0.20\textwidth} >{\raggedright\arraybackslash}X >{\raggedright\arraybackslash}X}
\toprule
 & \textbf{Posición protectoria} & \textbf{Posición de prevención del abuso} \\
\midrule
Idea central &
El sistema debe ofrecer una salida real a las personas humanas que, por circunstancias económicas o personales, no pueden afrontar regularmente sus obligaciones. &
La tutela del deudor no puede realizarse a costa de desconocer los derechos de los acreedores. \\
Principal preocupación &
Que el procedimiento concursal se vuelva inaccesible para quienes más lo necesitan. &
Que la flexibilización genere incentivos para incumplir obligaciones o utilizar el proceso con una finalidad distinta de la reorganización patrimonial. \\
Riesgo de una regulación unilateral &
Una regulación que solo priorice la satisfacción del crédito puede dejar sin respuesta a personas con una insolvencia genuina. &
Una regulación que solo contemple la protección del deudor podría perjudicar injustificadamente a los acreedores. \\
\bottomrule
\end{tabularx}
\end{table}

Ambas posiciones identifican problemas reales. Una regulación que solo contemple la protección del deudor podría perjudicar injustificadamente a los acreedores, pero un sistema que solo priorice la satisfacción del crédito puede dejar sin respuesta a personas que atraviesan una insolvencia genuina y que no cuentan con posibilidades materiales de reorganizar sus obligaciones.

\section{Posición propia}

La solución debe ubicarse en una posición intermedia. La persona humana no empresaria debe contar con un mecanismo efectivo para afrontar una situación de insolvencia, pero ese mecanismo debe estar condicionado a la transparencia patrimonial, a la buena fe y al cumplimiento de obligaciones razonables.

No se comparte la idea de que el concurso preventivo deba ser restringido por el solo hecho de que el deudor no desarrolle una actividad empresarial. La insolvencia puede afectar a cualquier persona y el artículo 2 de la LCQ reconoce expresamente esa posibilidad. Tampoco se considera adecuado que el sobreendeudamiento genere automáticamente un derecho a la cancelación de las obligaciones. El sistema debe distinguir entre la insolvencia provocada por contingencias económicas o personales y las situaciones en las que existen conductas fraudulentas o abusivas.

La postura sostenida es que \textbf{la buena fe debe ser un elemento central para diseñar una respuesta diferenciada}, pero no debe utilizarse como una fórmula abstracta. Es necesario establecer criterios objetivos que permitan examinar:
\begin{itemize}
    \item las causas del endeudamiento;
    \item la información suministrada por el deudor;
    \item la evolución de sus ingresos;
    \item la existencia de bienes no denunciados;
    \item la conducta de los acreedores.
\end{itemize}

En definitiva, el derecho concursal debe proteger al deudor honesto sin convertir esa protección en una ventaja injustificada frente a quienes le concedieron crédito.

\section*{Cuadro Cornell}
\addcontentsline{toc}{section}{Cuadro Cornell}
\begin{longtable}{QR}
\toprule
Preguntas / palabras clave & \textbf{Notas / respuestas} \\
\midrule
\endhead
¿Qué sostiene la posición protectoria? &
Que el derecho concursal debe dar una respuesta adecuada a la insolvencia de las personas humanas y no limitarse a reproducir las soluciones pensadas para las empresas. El sobreendeudamiento puede comprometer la subsistencia del deudor y su grupo familiar y provocar exclusión económica; y la respuesta jurídica no debería ser idéntica en todos los casos. \\[0.8em]
¿Qué es el crédito responsable? &
La idea de que el otorgamiento de financiación debe considerar la capacidad real de repago y evitar prácticas que favorezcan la acumulación de obligaciones difícilmente asumibles. No implica perdonar toda deuda de consumo ni que el acreedor pierda su derecho a cobrar. \\[0.8em]
¿Qué advierte la posición de prevención del abuso? &
Los riesgos de flexibilizar excesivamente el régimen: que el concurso permita liberarse mediante propuestas irrazonables o paralizar acciones individuales de cobro, debilitando la confianza en el crédito y la igualdad entre acreedores. \\[0.8em]
¿Cómo puede manifestarse el abuso? &
Ocultamiento de bienes, simulación de deudas, nuevos préstamos sin intención de pago, uso dilatorio del concurso y propuestas que no guardan relación razonable con las posibilidades económicas del deudor. \\[0.8em]
¿Cuál es la diferencia entre ambas posiciones? &
La protectoria teme que el procedimiento sea inaccesible para quienes más lo necesitan; la de prevención del abuso teme que la flexibilización incentive el incumplimiento o el uso del proceso con fines distintos de la reorganización. \\[0.8em]
¿Cuál es la posición propia? &
Intermedia: mecanismo efectivo condicionado a transparencia patrimonial, buena fe y cumplimiento de obligaciones razonables; la buena fe como elemento central, evaluada con criterios objetivos (causas del endeudamiento, información suministrada, evolución de ingresos, bienes no denunciados y conducta de los acreedores). \\
\midrule
\multicolumn{2}{p{\dimexpr0.92\textwidth+2\tabcolsep\relax}}{\textbf{Resumen:} La doctrina presenta una posición protectoria, que exige una respuesta específica para la persona humana sobreendeudada, y otra de prevención del abuso, que prioriza la tutela del crédito y la igualdad de los acreedores. Ambas identifican problemas reales; la posición propia es intermedia: proteger al deudor honesto, con buena fe evaluada mediante criterios objetivos, sin generar una ventaja injustificada frente a los acreedores.} \\
\bottomrule
\end{longtable}

% -------------------- D --------------------
\chapter[Jurisprudencia]{Jurisprudencia}

\section[Acceso al concurso: precedente de la Cámara de Azul]{El acceso al concurso de la persona humana sobreendeudada: el precedente de la Cámara de Apelación Civil y Comercial de Azul}

Entre los antecedentes mencionados en las versiones preliminares se encuentra el caso \textbf{``C., A. B. y otro/a s/ Concurso preventivo (pequeño)''}, atribuido a la Cámara de Apelación en lo Civil y Comercial de Azul, Sala II, con fecha 19 de junio de 2014.

Según la reconstrucción del caso que aparece en la referencia doctrinaria atribuida a Rossi (2019), se trataba de un matrimonio de jubilados que atravesaba una situación de sobreendeudamiento. El juzgado de primera instancia había rechazado la presentación concursal por considerar que no se encontraban cumplidos los requisitos formales del artículo 11 de la LCQ. La Cámara habría revocado esa decisión por considerarla prematura y habría otorgado un plazo para que los deudores completaran la documentación necesaria.

El interés del precedente reside en que permite observar la tensión entre el cumplimiento de las exigencias formales de la LCQ y la necesidad de brindar una respuesta judicial a personas humanas que atraviesan una situación de insolvencia.

La cuestión no consiste en eliminar los requisitos del artículo 11. Esos requisitos permiten conocer la situación patrimonial del deudor y proteger los derechos de los acreedores. El problema aparece cuando se los aplica sin considerar la naturaleza de la actividad del solicitante y la documentación que razonablemente puede exigírsele.

Desde esta perspectiva, una solución que permita completar la documentación antes de rechazar definitivamente la presentación puede resultar adecuada cuando los defectos son subsanables y no existe una conducta de ocultamiento o fraude.

% TODO: contenido incompleto o ambiguo en las notas originales (la referencia al caso Azul debe cotejarse con la sentencia original; Rossi sin datos completos)
\begin{important}
\textbf{Aclaración metodológica:} la referencia al caso se encuentra reconstruida a partir de una fuente doctrinaria secundaria y no de la sentencia completa. Por ello, antes de presentar el trabajo como definitivo, corresponde obtener y cotejar el texto original para confirmar la carátula, la sala, la fecha, los fundamentos y el alcance exacto de lo resuelto. No corresponde atribuir al tribunal argumentos que no puedan comprobarse en la resolución.
\end{important}

\section[Control judicial de los acuerdos preventivos]{El control judicial de los acuerdos preventivos: ``Arcángel Maggio S.A.'' y ``Sociedad Comercial del Plata S.A.''}

La problemática del abuso no se limita al acceso al concurso. También aparece en la etapa de negociación y homologación del acuerdo preventivo.

En este punto resultan relevantes los precedentes de la Corte Suprema de Justicia de la Nación \textbf{``Arcángel Maggio S.A. s/ concurso preventivo s/ incidente de impugnación al acuerdo preventivo''}, Fallos 330:834, y \textbf{``Sociedad Comercial del Plata S.A. y otros s/ concurso preventivo''}, Fallos 332:2339.

Estos precedentes son relevantes para estudiar el alcance del control judicial de las propuestas y acuerdos preventivos. Su inclusión no significa que hayan resuelto directamente un caso de sobreendeudamiento de consumidores o de personas humanas no empresarias. Su utilidad es otra: permiten analizar los límites que la legislación impone a la homologación de acuerdos que puedan resultar abusivos o contrarios a la ley.

El artículo 52, inciso 4, de la LCQ establece expresamente que el juez no homologará una propuesta abusiva o en fraude a la ley. Por ello, la obtención de las mayorías legales no elimina el control judicial. Esta facultad adquiere especial importancia cuando se analiza la posibilidad de que una persona humana presente una propuesta de quita o espera que no guarde relación con sus posibilidades reales de pago o que, por sus características, implique un sacrificio injustificado para los acreedores.

La solución no puede consistir en rechazar toda propuesta que contenga una quita importante. En muchos casos, una propuesta de pago parcial puede ser más conveniente para los acreedores que una quiebra sin activos realizables. Pero tampoco puede aceptarse cualquier propuesta por el solo hecho de haber obtenido las mayorías legales.

El juez debe examinar las circunstancias del caso, la situación patrimonial del deudor, la razonabilidad de la propuesta y la eventual existencia de fraude o abuso.

\section{Valoración conjunta de la jurisprudencia}

Los antecedentes analizados permiten identificar dos dimensiones del problema.

Por un lado, el caso atribuido a la Cámara de Azul pone de manifiesto la necesidad de evitar que los requisitos formales se conviertan en obstáculos desproporcionados para el acceso de las personas humanas al proceso concursal.

Por otro lado, los precedentes de la Corte Suprema sobre control de acuerdos preventivos permiten advertir que la protección del deudor no puede traducirse en una homologación automática de cualquier propuesta.

Se coincide con una interpretación que combine ambos aspectos. El juez debe evitar que el formalismo impida injustificadamente el acceso al procedimiento, pero también debe controlar que el deudor no utilice el concurso para perjudicar a sus acreedores.

La finalidad del proceso concursal no es favorecer a una de las partes de manera indiscriminada, sino ofrecer una solución colectiva, ordenada y jurídicamente controlada a una situación de insolvencia.

\section*{Cuadro Cornell}
\addcontentsline{toc}{section}{Cuadro Cornell}
\begin{longtable}{QR}
\toprule
Preguntas / palabras clave & \textbf{Notas / respuestas} \\
\midrule
\endhead
¿Qué caso se atribuye a la Cámara de Azul y qué se reconstruye de él? &
``C., A. B. y otro/a s/ Concurso preventivo (pequeño)'', Cámara de Apelación en lo Civil y Comercial de Azul, Sala II, 19 de junio de 2014. Según la referencia atribuida a Rossi (2019), un matrimonio de jubilados sobreendeudado vio rechazada su presentación en primera instancia por incumplir el art.\ 11; la Cámara habría revocado por considerar la decisión prematura y habría otorgado plazo para completar la documentación. \\[0.8em]
¿Qué cautela metodológica corresponde respecto de este caso? &
La referencia proviene de una fuente doctrinaria secundaria y no de la sentencia completa; debe cotejarse el texto original (carátula, sala, fecha, fundamentos y alcance) y no atribuir al tribunal argumentos no comprobables. \\[0.8em]
¿Qué tensión muestra el precedente? &
La existente entre las exigencias formales del art.\ 11 y la necesidad de dar respuesta judicial a personas humanas insolventes. Permitir completar la documentación antes del rechazo definitivo es adecuado cuando los defectos son subsanables y no hay ocultamiento o fraude. \\[0.8em]
¿Por qué se incluyen ``Arcángel Maggio S.A.'' y ``Sociedad Comercial del Plata S.A.''? &
Fallos 330:834 y 332:2339 de la CSJN permiten estudiar el alcance del control judicial de propuestas y acuerdos preventivos y los límites a la homologación de acuerdos abusivos o contrarios a la ley; no resolvieron directamente casos de sobreendeudamiento de personas humanas no empresarias. \\[0.8em]
¿Qué debe examinar el juez ante una propuesta? &
Las circunstancias del caso, la situación patrimonial del deudor, la razonabilidad de la propuesta y la eventual existencia de fraude o abuso. No debe rechazarse toda quita importante (un pago parcial puede convenir más que una quiebra sin activos realizables), ni aceptarse cualquier propuesta solo por tener mayorías. \\
\midrule
\multicolumn{2}{p{\dimexpr0.92\textwidth+2\tabcolsep\relax}}{\textbf{Resumen:} La jurisprudencia analizada muestra dos dimensiones: evitar que el formalismo obstaculice desproporcionadamente el acceso de las personas humanas al concurso (caso de Azul, a cotejar con la sentencia original) y controlar que la protección del deudor no implique homologar cualquier propuesta (precedentes de la CSJN). La finalidad del proceso es una solución colectiva, ordenada y controlada.} \\
\bottomrule
\end{longtable}

% -------------------- E --------------------
\chapter[Análisis crítico]{Análisis crítico}

\section{La diferencia entre la insolvencia empresarial y la insolvencia de la persona humana}

Uno de los principales problemas del régimen actual es que somete a sujetos con realidades económicas muy diferentes a una estructura procedimental común.

Una empresa puede tener trabajadores, proveedores, instalaciones, activos productivos, contratos en ejecución y una organización contable. La persona humana no empresaria, en cambio, puede tener como único ingreso un salario o una jubilación y haber contraído deudas para afrontar gastos familiares.

En la empresa, la reorganización puede tener como objetivo preservar una unidad productiva y mantener su funcionamiento. En la persona humana no empresaria, el problema suele ser distinto: ordenar las deudas, proteger los ingresos necesarios para la subsistencia y establecer una forma realista de afrontar las obligaciones.

Esto no significa que el régimen común sea necesariamente inaplicable. Significa que sus reglas deben interpretarse y, en su caso, reformarse teniendo en cuenta las diferencias entre ambos sujetos.

\section{Las dificultades económicas y procesales de acceso}

El procedimiento concursal exige gastos y actividades que pueden resultar difíciles de afrontar para una persona que ya se encuentra en cesación de pagos. Entre ellos se encuentran los gastos de publicidad, la asistencia profesional, la preparación de documentación, la intervención de la sindicatura y los costos vinculados con la verificación de créditos.

El problema es especialmente visible cuando el deudor carece de bienes realizables y depende exclusivamente de ingresos periódicos para sostener sus necesidades básicas. La finalidad del proceso es ordenar la situación patrimonial, pero sus costos pueden convertirse en un obstáculo para llegar a esa solución.

No se sostiene que todos los costos deban eliminarse. La intervención profesional y la publicidad cumplen funciones necesarias. Sin embargo, se considera que el legislador debería prever mecanismos diferenciados para los casos de insolvencia de personas humanas no empresarias, con procedimientos más sencillos y gastos proporcionales a la capacidad económica del solicitante.

\section{El riesgo de abuso y la protección de los acreedores}

La necesidad de facilitar el acceso al concurso no debe llevar a desconocer los derechos de los acreedores.

La apertura del proceso produce efectos relevantes: permite organizar colectivamente el tratamiento de las obligaciones y genera consecuencias sobre las acciones individuales. Por eso, puede existir un incentivo para que un deudor utilice el concurso con una finalidad meramente dilatoria. También puede existir abuso en la presentación de información patrimonial incompleta, en el ocultamiento de bienes, en la simulación de obligaciones o en la formulación de propuestas que no respeten los límites legales.

La prevención de estas conductas exige controles eficaces. La sindicatura debe poder examinar la situación patrimonial y las causas de la insolvencia, mientras que el juez debe ejercer las facultades que le otorga la LCQ para evitar el fraude y la homologación de acuerdos abusivos.

La solución no consiste en presumir la mala fe de toda persona sobreendeudada. Tampoco consiste en considerar que todo deudor es necesariamente una víctima del sistema financiero. \textbf{La buena fe debe examinarse en cada caso, a partir de circunstancias concretas.}

\section{El crédito responsable y la conducta de los proveedores financieros}

El análisis del sobreendeudamiento no puede limitarse a la conducta del deudor.

En las relaciones de consumo, los proveedores profesionales de crédito cuentan con herramientas de información y evaluación que no necesariamente posee el consumidor. Pueden analizar ingresos, antecedentes crediticios, obligaciones existentes y capacidad estimada de repago. La expansión de los préstamos digitales y de las ofertas de financiación inmediata plantea nuevos interrogantes sobre la evaluación de la solvencia del consumidor.

El otorgamiento sucesivo de créditos sin considerar la capacidad real de pago puede contribuir a una situación de sobreendeudamiento. Sin embargo, no corresponde afirmar que todo crédito concedido a una persona que luego incumple sea irresponsable ni que exista automáticamente responsabilidad del prestamista.

El debate sobre el crédito responsable exige examinar los deberes legales de información, transparencia y trato digno, así como la conducta concreta del proveedor y las circunstancias de cada operación.

Se considera que una futura regulación debería fortalecer los mecanismos de evaluación de solvencia y de prevención del endeudamiento excesivo, sin impedir injustificadamente el acceso al crédito.

\section{La ausencia de un régimen específico de segunda oportunidad}

Una de las principales diferencias entre el sistema argentino y algunos modelos extranjeros es la falta de un procedimiento general de exoneración del pasivo insatisfecho equivalente al existente en España.

El régimen argentino contempla la conclusión de los procesos concursales y los efectos de la quiebra, pero no ofrece, en los mismos términos que la EPI española, un mecanismo general de liberación del pasivo remanente para la persona humana de buena fe que ha atravesado una insolvencia y cumplido las condiciones legales.

Esto puede generar una situación de exclusión económica prolongada. Una persona puede haber liquidado sus bienes o haber realizado esfuerzos razonables para pagar, pero continuar condicionada por obligaciones que no tiene posibilidades materiales de satisfacer.

El derecho comparado permite considerar una alternativa: establecer un mecanismo de segunda oportunidad sujeto a requisitos estrictos, que contemple la buena fe del deudor, la transparencia patrimonial, la colaboración con el proceso y el cumplimiento de un plan de pagos cuando corresponda.

No se trataría de premiar el incumplimiento, sino de reconocer que existen situaciones en las que mantener indefinidamente una deuda impagable no beneficia ni al deudor ni a los acreedores.

\section{Propuestas para mejorar el sistema argentino}

A partir del análisis realizado, se considera que una reforma legislativa debería contemplar los siguientes aspectos:

\begin{enumerate}[label=\alph*)]
    \item \textbf{Un procedimiento específico para personas humanas no empresarias.} Debería crearse un régimen diferenciado para quienes atraviesan una situación de insolvencia vinculada con obligaciones personales o de consumo, sin actividad empresarial organizada.
    \item \textbf{Requisitos documentales adecuados a la realidad del deudor.} La información patrimonial debe ser completa y verificable, pero no debería exigirse documentación contable que no corresponda a la actividad del solicitante.
    \item \textbf{Publicidad accesible y digital.} Deberían fortalecerse los medios oficiales de publicidad electrónica y los mecanismos de notificación a los acreedores conocidos, preservando la transparencia del proceso.
    \item \textbf{Una instancia previa de negociación.} Podría incorporarse un procedimiento de conciliación o renegociación que permita alcanzar acuerdos antes de iniciar un proceso judicial más complejo, siempre que se asegure la participación de los acreedores y no se utilice como mecanismo dilatorio.
    \item \textbf{Controles objetivos de buena fe.} La ley debería establecer criterios para evaluar la conducta del deudor, incluyendo la información patrimonial, las causas del endeudamiento, la existencia de actos de ocultamiento y la colaboración durante el procedimiento.
    \item \textbf{Un mecanismo de segunda oportunidad.} Debería estudiarse la incorporación de una exoneración del pasivo insatisfecho para personas humanas de buena fe, bajo condiciones legales claras y con exclusiones para las obligaciones que el legislador considere que deben recibir un tratamiento especial.
    \item \textbf{Prevención del crédito irresponsable.} La regulación debería promover una evaluación razonable de la capacidad de pago del consumidor y mecanismos de información que permitan evitar la acumulación descontrolada de obligaciones.
\end{enumerate}

Estas propuestas no buscan eliminar la tutela del crédito. Por el contrario, pretenden construir un sistema más equilibrado, en el que el deudor tenga una posibilidad real de superar la insolvencia y los acreedores cuenten con garantías suficientes para evitar maniobras abusivas.

\section*{Cuadro Cornell}
\addcontentsline{toc}{section}{Cuadro Cornell}
\begin{longtable}{QR}
\toprule
Preguntas / palabras clave & \textbf{Notas / respuestas} \\
\midrule
\endhead
¿En qué se diferencian la insolvencia empresarial y la de la persona humana? &
La empresa puede tener trabajadores, proveedores, instalaciones, activos productivos, contratos y organización contable, y la reorganización puede buscar preservar una unidad productiva. La persona humana no empresaria puede tener como único ingreso un salario o jubilación; su problema suele ser ordenar deudas, proteger los ingresos necesarios para la subsistencia y afrontar las obligaciones de forma realista. \\[0.8em]
¿Qué obstáculos económicos y procesales existen? &
Gastos de publicidad, asistencia profesional, preparación de documentación, intervención de la sindicatura y costos de verificación de créditos, especialmente graves si el deudor carece de bienes realizables. Se propone prever mecanismos diferenciados, más sencillos y con gastos proporcionales a la capacidad económica, sin eliminar todos los costos. \\[0.8em]
¿Cómo equilibrar acceso y protección de los acreedores? &
Mediante controles eficaces (sindicatura y juez) contra la información incompleta, el ocultamiento de bienes, la simulación de obligaciones y las propuestas fuera de los límites legales, sin presumir la mala fe de todo sobreendeudado ni considerar víctima a todo deudor: la buena fe se examina caso por caso. \\[0.8em]
¿Qué papel juega la conducta de los proveedores de crédito? &
Cuentan con herramientas para evaluar la solvencia, y el otorgamiento sucesivo de créditos sin considerar la capacidad de pago puede contribuir al sobreendeudamiento; pero no todo crédito a quien luego incumple es irresponsable ni genera automáticamente responsabilidad del prestamista. Debe examinarse deberes de información, transparencia y trato digno, y la conducta concreta. \\[0.8em]
¿Qué carencia tiene el sistema argentino frente al español? &
La falta de un procedimiento general de exoneración del pasivo insatisfecho equivalente a la EPI, lo que puede generar exclusión económica prolongada. Una alternativa sería un mecanismo de segunda oportunidad con requisitos estrictos: buena fe, transparencia patrimonial, colaboración y plan de pagos cuando corresponda. \\[0.8em]
¿Qué reformas se proponen? &
(a) procedimiento específico para personas humanas no empresarias; (b) requisitos documentales adecuados; (c) publicidad accesible y digital; (d) instancia previa de negociación; (e) controles objetivos de buena fe; (f) mecanismo de segunda oportunidad; (g) prevención del crédito irresponsable. \\
\midrule
\multicolumn{2}{p{\dimexpr0.92\textwidth+2\tabcolsep\relax}}{\textbf{Resumen:} El régimen somete a sujetos muy distintos a una estructura común, con costos y exigencias que pueden obstaculizar el acceso del deudor de buena fe, y sin segunda oportunidad. Se propone una reforma equilibrada que combine acceso proporcionado, controles objetivos de buena fe, crédito responsable y exoneración del pasivo remanente, sin eliminar la tutela del crédito.} \\
\bottomrule
\end{longtable}

% -------------------- F --------------------
\chapter[Conclusiones]{Conclusiones}

La pregunta que orientó esta investigación fue si el concurso preventivo regulado por la Ley 24.522 constituye una herramienta adecuada para proteger a las personas humanas no empresarias afectadas por el sobreendeudamiento, sin desatender los derechos de los acreedores ni facilitar el abuso del proceso.

Luego de analizar el régimen normativo, las posiciones doctrinarias, los antecedentes jurisprudenciales y algunas soluciones del derecho comparado, se considera que \textbf{el sistema argentino ofrece una respuesta insuficiente para las particularidades de la insolvencia de las personas humanas no empresarias}.

El artículo 2 de la LCQ reconoce ampliamente la posibilidad de que una persona humana acceda al concurso preventivo. Sin embargo, esa apertura subjetiva no fue acompañada por un procedimiento integralmente adaptado a quienes no desarrollan una actividad empresarial.

Los pequeños concursos de los artículos 288 y 289 representan una simplificación, pero no resuelven el problema de fondo. El procedimiento continúa conservando etapas y exigencias que pueden resultar desproporcionadas para un trabajador, jubilado o consumidor cuyo patrimonio se encuentra compuesto principalmente por ingresos personales y obligaciones de consumo.

Por eso, el régimen actual puede convertirse en una herramienta de difícil acceso para el deudor de buena fe. Los costos del procedimiento, las exigencias documentales y las cargas de publicidad pueden obstaculizar la posibilidad de obtener una solución colectiva.

Al mismo tiempo, la protección de la persona humana no puede justificar cualquier forma de liberación de obligaciones. Los acreedores tienen derechos que deben ser respetados y el proceso concursal no debe transformarse en un instrumento para paralizar ejecuciones, ocultar bienes o imponer acuerdos abusivos.

La solución debe construirse a partir de un equilibrio:
\begin{itemize}
    \item El deudor debe actuar con transparencia, informar su situación patrimonial y colaborar con el procedimiento.
    \item Los acreedores deben tener la posibilidad de verificar sus créditos, participar en la negociación y cuestionar las propuestas contrarias a la ley.
    \item El juez debe garantizar que el proceso cumpla su finalidad y evitar tanto el formalismo excesivo como el abuso del derecho.
\end{itemize}

El derecho comparado, especialmente la experiencia española con la Exoneración del Pasivo Insatisfecho, demuestra que es posible pensar mecanismos de segunda oportunidad sujetos a requisitos y controles. Esa experiencia no debe trasladarse automáticamente, pero constituye una referencia para debatir una reforma del sistema argentino.

\begin{important}
\textbf{Conclusión:} el concurso preventivo de la persona humana no debe ser prohibido ni convertido en un privilegio, sino complementado con un régimen específico de tratamiento del sobreendeudamiento.
\end{important}

Ese régimen debería ofrecer procedimientos accesibles, requisitos adecuados a la realidad económica del deudor, instancias de negociación, controles de buena fe y, cuando corresponda, mecanismos de exoneración del pasivo remanente.

La finalidad no debe ser liberar indiscriminadamente a las personas de sus obligaciones, sino permitir que quien atravesó una situación genuina de insolvencia pueda recuperar su capacidad económica, sin desconocer los derechos de quienes le otorgaron crédito.

En definitiva, el desafío del derecho concursal contemporáneo consiste en superar una visión exclusivamente patrimonial de la insolvencia y reconocer que detrás de determinados procesos concursales existe una persona humana cuya situación económica puede comprometer su subsistencia y su reinserción social. Pero esa humanización del proceso debe desarrollarse dentro de un marco de legalidad, buena fe y respeto por los derechos de todos los involucrados.

\section*{Cuadro Cornell}
\addcontentsline{toc}{section}{Cuadro Cornell}
\begin{longtable}{QR}
\toprule
Preguntas / palabras clave & \textbf{Notas / respuestas} \\
\midrule
\endhead
¿Cuál es la conclusión general sobre el sistema argentino? &
Ofrece una respuesta insuficiente para las particularidades de la insolvencia de las personas humanas no empresarias: el art.\ 2 abre el acceso, pero no hay un procedimiento integralmente adaptado; los pequeños concursos simplifican pero no resuelven el problema de fondo. \\[0.8em]
¿Por qué puede ser de difícil acceso para el deudor de buena fe? &
Por los costos del procedimiento, las exigencias documentales y las cargas de publicidad. \\[0.8em]
¿Qué límites tiene la protección de la persona humana? &
No justifica cualquier liberación de obligaciones; el concurso no debe servir para paralizar ejecuciones, ocultar bienes o imponer acuerdos abusivos. \\[0.8em]
¿Cuál es el equilibrio propuesto entre deudor, acreedores y juez? &
El deudor actúa con transparencia y colabora; los acreedores verifican créditos, participan en la negociación y cuestionan propuestas contrarias a la ley; el juez evita tanto el formalismo excesivo como el abuso del derecho. \\[0.8em]
¿Qué solución se propone? &
Que el concurso de la persona humana no sea prohibido ni un privilegio, sino complementado con un régimen específico con procedimientos accesibles, requisitos adecuados, instancias de negociación, controles de buena fe y, cuando corresponda, exoneración del pasivo remanente. \\
\midrule
\multicolumn{2}{p{\dimexpr0.92\textwidth+2\tabcolsep\relax}}{\textbf{Resumen:} El sistema actual es insuficiente para la persona humana no empresaria. La solución es un régimen complementario específico, con equilibrio entre la recuperación del deudor honesto y los derechos de los acreedores, dentro de un marco de legalidad y buena fe.} \\
\bottomrule
\end{longtable}

% ==================== BIBLIOGRAFÍA ====================
\chapter*{Bibliografía y fuentes}
\addcontentsline{toc}{chapter}{Bibliografía y fuentes}
\markboth{Bibliografía y fuentes}{Bibliografía y fuentes}

\section*{Legislación}
\begin{itemize}
    \item Constitución de la Nación Argentina, especialmente artículos 16, 18, 42 y 75 inciso 22.
    \item Ley 24.522 de Concursos y Quiebras y sus modificatorias.
    \item Ley 24.240 de Defensa del Consumidor y sus modificatorias.
    \item Ley 26.994, Código Civil y Comercial de la Nación, especialmente artículos 9 y 10.
    \item Ley 27.170, modificatoria del régimen de pequeños concursos.
    \item Reino de España, Ley 16/2022, de 5 de septiembre, de reforma del texto refundido de la Ley Concursal.
    \item Unión Europea, Directiva (UE) 2019/1023.
\end{itemize}

\section*{Doctrina}
% TODO: contenido incompleto o ambiguo en las notas originales (referencias bibliográficas a completar y cotejar)
\begin{itemize}
    \item Junyent Bas, Francisco, y Molina Sandoval, Carlos A. \textit{Ley de Concursos y Quiebras comentada}. Referencia bibliográfica a completar y cotejar según la edición efectivamente consultada.
    \item Rivera, Julio César. \textit{Instituciones de Derecho Concursal}. Referencia bibliográfica a completar según la edición y los tomos efectivamente consultados.
    \item Rouillon, Adolfo A. N. \textit{Régimen de Concursos y Quiebras. Ley 24.522}. Referencia bibliográfica a completar según la edición efectivamente consultada.
    \item Rossi, [nombre completo y datos editoriales pendientes]. Trabajo doctrinario citado en las versiones preliminares para la reconstrucción del precedente de la Cámara de Apelación Civil y Comercial de Azul. Completar título, publicación, año y páginas.
    \item Apuntes y transcripciones de clase de la asignatura Concursos y Quiebras, Facultad de Derecho, Universidad de Buenos Aires, segundo bimestre de 2026.
\end{itemize}

\section*{Jurisprudencia}
\begin{itemize}
    \item Cámara de Apelación en lo Civil y Comercial de Azul, Sala II, ``C., A. B. y otro/a s/ Concurso preventivo (pequeño)'', 19 de junio de 2014. \textbf{Referencia pendiente de cotejo con la sentencia original y de completar los datos de publicación.}
    \item Corte Suprema de Justicia de la Nación, ``Arcángel Maggio S.A. s/ concurso preventivo s/ incidente de impugnación al acuerdo preventivo'', Fallos 330:834.
    \item Corte Suprema de Justicia de la Nación, ``Sociedad Comercial del Plata S.A. y otros s/ concurso preventivo'', Fallos 332:2339.
\end{itemize}

\section*{Declaración sobre el uso de inteligencia artificial}
En la elaboración del presente trabajo se utilizaron herramientas de inteligencia artificial como asistencia para organizar los materiales de estudio, comparar versiones preliminares, identificar problemas jurídicos, revisar la redacción y estructurar el análisis normativo, doctrinario y jurisprudencial. La selección definitiva de las fuentes, la verificación de las citas y la responsabilidad por las opiniones y conclusiones corresponden a los autores del trabajo.

\end{document}
